\section{Automorphic support and local presentations}
\label{sec:sr-support-propagation}
The arithmetic endpoint requires one automorphic support to propagate
through a connected deformation space. We first give the specialization
and connectedness argument, then construct a single patched module
from explicit finite-level inputs, and finally prove the determinant
algebra used at the paired locus. These proofs do not replace the
construction of the finite-level automorphic modules or the Galois
presentations to which they apply.

\subsection{The characteristic-prime propagation mechanism}
\label{sec:sr-characteristic-propagation}
We isolate and prove the prime-selection and connectedness arguments in
the rank-three endpoint. These arguments do not invoke an automorphy
lifting theorem. Their application still requires the specific
deformation presentations and the Hecke-support implications stated
at the end of this section.

\subsubsection{A one-dimensional characteristic-prime specialization}
\begin{srlemma}\label{lem:sr-nice-prime}
Let $(R,\mathfrak m)$ be a Noetherian local ring, let
$\varpi\in\mathfrak m$, and let $D\ge1$.
Suppose that $I,\mathfrak a_1,\ldots,\mathfrak a_s$ are ideals with
\[
 \dim R/I\ge D+1,\qquad
                  \dim R/\mathfrak a_i<D\quad(1\le i\le s).
\]
There is a prime $\mathfrak p$ such that
\[
 I+(\varpi)\subseteq\mathfrak p,\qquad
 \dim R/\mathfrak p=1,\qquad
                 \mathfrak a_i\not\subseteq\mathfrak p
                              \quad(1\le i\le s).
\]
\end{srlemma}
\begin{proof}
The local dimension inequality for one equation gives
$\dim R/(I,\varpi)\ge D$.
One proof of that inequality is to lift a system of parameters of
$R/(I,\varpi)$ to $R/I$ and adjoin $\varpi$; the resulting maximal-primary
ideal has at most one additional generator, and Krull's height theorem
bounds the dimension. Since a Noetherian ring has finitely many minimal
primes, choose a minimal prime $\mathfrak q$ of $(I,\varpi)$ with
$\dim R/\mathfrak q\ge D$.
Then $\mathfrak a_i\not\subseteq\mathfrak q$ for every $i$.
Choose $x_i\in\mathfrak a_i\setminus\mathfrak q$ and
$g\in\mathfrak m\setminus\mathfrak q$, and put $f=g\prod_i x_i$.
Thus $f\in\mathfrak m\setminus\mathfrak q$.

Choose a prime $\mathfrak p$ maximal among the primes containing
$\mathfrak q$ and avoiding $f$; the ascending-chain condition supplies
a maximal element. Put $A=R/\mathfrak p$, with maximal ideal
$\mathfrak n$ and image $\bar f$ of $f$.
This is a local domain, $0\ne\bar f\in\mathfrak n$, and every nonzero
prime of $A$ contains $\bar f$, by maximality of $\mathfrak p$.
We claim $\dim A=1$.

For any $0\ne x\in\mathfrak n$, choose a prime $P$ minimal over $(x)$.
It is nonzero and has height at most one by the principal ideal theorem.
It contains $\bar f$. Choose $Q$ minimal over $(\bar f)$ inside $P$.
Since $\bar f\ne0$, $Q\ne0$. A strict inclusion $Q\subsetneq P$ would
give the chain $0\subsetneq Q\subsetneq P$, contrary to the height
bound for $P$. Hence $P=Q$.
It follows that every element of $\mathfrak n$ lies in the union of
the finitely many primes minimal over $(\bar f)$.
Finite prime avoidance places $\mathfrak n$ inside one of them, so
$\mathfrak n$ itself has height at most one.
The nonzero element $\bar f\in\mathfrak n$ gives
$0\subsetneq\mathfrak n$, so the dimension is exactly one.
Finally $f\notin\mathfrak p$ forces
$x_i\notin\mathfrak p$ for every $i$, proving the required avoidance.
\end{proof}
For $\varpi=\ell$, this produces a characteristic-$\ell$ prime.
The conclusion concerns a one-dimensional quotient; it does not assert
that the selected prime has height one.
The source proof is
\href{https://github.com/ChristianKudeba/alpoge-rh-verification/blob/alpoge-source-7cbf7a0c/src/Development/AutomorphicArithmetic/AutomorphicVarieties/CurveDVRSpecialization/NicePrimeEx.lean}
{\srcid{CurveDVRSpecialization.NicePrimeEx}}.

\subsubsection{Local propagation across ordinary and exceptional components}
For a Noetherian ring $R$ and a prime $P$, write
$d(P)=\dim R/P$. Fix an integer $D$, an ideal $\mathfrak s$, and five
sets of primes $\mathcal T,\mathcal E,\mathcal G,\mathcal N,\mathcal W$.
Think of these as the ordinary locus, exceptional locus, supported
locus, selected specialization primes, and thick exceptional components.
The following lemma states every hypothesis used by the local argument.

Let $\mathcal C=\operatorname{Min}(\mathfrak s)$.
Two members $P,Q\in\mathcal C$ are adjacent if
$\dim R/(P+Q)\ge D$.
Two selected exceptional primes $p,q\in\mathcal N\cap\mathcal E$ are
linked if there is $W\in\mathcal C\cap\mathcal W$ with $W\subseteq p,q$.
A chain in either relation means a finite chain, allowing length zero.

\begin{srlemma}\label{lem:sr-local-propagation}
Assume:
\begin{enumerate}
\item $\mathcal G$ is closed under specialization, and $\mathcal T$
is closed under generalization.
\item Every $p\in\mathcal N$ belongs to $\mathcal T\cup\mathcal E$.
\item If $I\in\mathcal G$ and $d(I)\ge D$, there is
$p\in\mathcal N$ containing $I$.
\item If $W\in\mathcal C\cap\mathcal W$, $W\subseteq I$, and
$d(I)\ge D$, there is $p\in\mathcal N\cap\mathcal E$ containing $I$.
\item Every minimal prime of $R$ belonging to $\mathcal T$ is in
$\mathcal C$. Every two members of $\mathcal C$ are connected by an
adjacency chain. Every member of $\mathcal C$ is in
$\mathcal T\cup\mathcal W$.
\item If $p\in\mathcal N\cap\mathcal T\cap\mathcal G$, every prime
$P\subseteq p$ belongs to $\mathcal G$.
\item Suppose $p,q\in\mathcal N\cap\mathcal E$ are linked by a chain,
and some $P\subseteq p$ belongs to $\mathcal T\cap\mathcal G$.
Then every $P'\subseteq q$ belonging to $\mathcal T$ belongs to
$\mathcal G$.
\end{enumerate}
If $I\in\mathcal T\cap\mathcal G$ and $d(I)\ge D$, then every minimal
prime of $R$ in $\mathcal T$ belongs to $\mathcal G$.
\end{srlemma}
\begin{proof}
First observe a consequence valid for every $p\in\mathcal N$:
if some $Q\subseteq p$ belongs to $\mathcal T\cap\mathcal G$,
then every $Q'\subseteq p$ in $\mathcal T$ belongs to $\mathcal G$.
If $p\in\mathcal T$, specialization gives $p\in\mathcal G$ and
hypothesis 6 applies. If $p\in\mathcal E$, hypothesis 7 applies to
the chain of length zero.

Choose a minimal prime $P_0$ of $R$ inside $I$.
Generalization gives $P_0\in\mathcal T$, hence $P_0\in\mathcal C$.
Hypothesis 3 supplies $p\in\mathcal N$ above $I$; the observation
with witness $I$ proves $P_0\in\mathcal G$.

We propagate the following information along an adjacency chain.
A component $Q\in\mathcal C$ is informed if either
$Q\in\mathcal T\cap\mathcal G$, or $Q\in\mathcal W$ and there are
selected exceptional primes $p_0,p$ joined by a link chain such that
some ordinary supported prime lies below $p_0$ and $Q\subseteq p$.
Informed ordinary components are supported, by definition.

Suppose $Q$ is informed and adjacent to $Q'$.
Choose a prime $M$ above $Q+Q'$ with $d(M)\ge D$; a minimal prime
of this quotient with maximal quotient dimension suffices.
If $Q\in\mathcal T\cap\mathcal G$, specialization makes $M$ supported.
Hypothesis 3 gives a selected $p\supseteq M$, and the first observation
shows that every ordinary prime below $p$ is supported.
If $Q$ is informed through an exceptional link chain, hypothesis 4
gives a selected exceptional $p\supseteq M$.
The old endpoint and $p$ are linked through $Q$; append this link to
the old chain and apply hypothesis 7.
Again every ordinary prime below $p$ is supported.
Thus if $Q'\in\mathcal T$, then $Q'\in\mathcal G$.

If $Q'\notin\mathcal T$, hypothesis 5 gives $Q'\in\mathcal W$.
Hypothesis 4 gives a selected exceptional $p'\supseteq M$.
When $Q$ is ordinary supported, $Q\subseteq p'$ is the initial
ordinary supported witness, so $Q'$ is informed with a chain of
length zero. When $Q$ is thick and informed, append the link from
the previous endpoint to $p'$ through $Q$.
This informs $Q'$ as well.

Starting at the supported ordinary component $P_0$, hypothesis 5
therefore informs every component in $\mathcal C$.
Every minimal prime of $R$ in $\mathcal T$ lies in $\mathcal C$
and is supported, which proves the conclusion.
\end{proof}
This is the entire argument of
\href{https://github.com/ChristianKudeba/alpoge-rh-verification/blob/alpoge-source-7cbf7a0c/src/Development/AutomorphicArithmetic/AutomorphicVarieties/CurveDVRSpecialization/LocProp.lean}
{\srcid{CurveDVRSpecialization.LocProp}}.
The separate propagation through thick exceptional primes is essential:
ordinary propagation at a single specialization does not cross those
components by itself.

\subsubsection{Global propagation from a connected characteristic-prime core}
\begin{srlemma}\label{lem:sr-global-propagation}
Let $X$ be a Noetherian topological space. Fix subsets
$T,G,C,D,E\subseteq X$.
Call an irreducible component relevant if it meets $T$.
Suppose:
\begin{enumerate}
\item $D\subseteq C$ and $C$ is connected.
\item Every closed subset meeting $C$ meets $D$.
\item $C\subseteq\overline{D\setminus E}$.
\item Every relevant irreducible component meets $C$.
\item If relevant components $Y,Y'$ meet at $x\in D$ and $Y\subseteq G$,
then $Y'\subseteq G$.
\item Every $x\in D\setminus E$ lies on a relevant component.
\item Some relevant component is contained in $G$.
\end{enumerate}
Then $T\subseteq G$.
\end{srlemma}
\begin{proof}
There are finitely many irreducible components.
Let $A$ be the union of all relevant components contained in $G$, and
let $B$ be the union of the other relevant components.
These are closed sets. Hypotheses 3 and 6 give
$C\subseteq A\cup B$.
The seed and hypothesis 4 imply $C\cap A\ne\varnothing$.
If a relevant component were not contained in $G$, hypothesis 4
would also give $C\cap B\ne\varnothing$.
Connectedness would then give $C\cap A\cap B\ne\varnothing$.
Hypothesis 2, applied to the closed set $A\cap B$, would give
$x\in D\cap A\cap B$.
Choose a component in each of the two unions containing $x$.
Hypothesis 5 would make the component from $B$ supported,
contradicting its definition.
Thus every relevant component is contained in $G$.
Every point of $T$ lies on a relevant component, proving the claim.
\end{proof}
This is
\href{https://github.com/ChristianKudeba/alpoge-rh-verification/blob/alpoge-source-7cbf7a0c/src/Development/AutomorphicArithmetic/AutomorphicVarieties/CurveDVRSpecialization/GlobProp.lean}
{\srcid{CurveDVRSpecialization.GlobProp}}.
The proof works with preconnected $C$; no nonemptiness is needed beyond
the seed and the hypotheses that the seed component meets $C$.

Here is how local charts give hypothesis 5. Let $T$ be open and $G$
closed. At $x\in D$, suppose that a flat local chart
$f:\operatorname{Spec}R\to X$ maps its closed point to $x$.
Its minimal primes map to the generic points of components through $x$,
and every such component is obtained this way.
For an affine neighborhood this follows by factoring through
$\mathcal O_{X,x}\to R$: the map is flat and local, hence faithfully
flat; lying over and going down give the two assertions.
Suppose every minimal prime of $R$ has quotient dimension at least
the threshold $D_0$ in Lemma~\ref{lem:sr-local-propagation}.
A relevant component $Y$ through $x$ has a minimal-prime lift $P$.
Its generic point is in the open set $T$, so $P\in f^{-1}T$.
If $Y\subseteq G$, then $P\in f^{-1}G$.
The local lemma now places every ordinary minimal-prime lift in
$f^{-1}G$. A lift of the generic point of any other relevant
component $Y'$ through $x$ is one of these primes.
Closedness of $G$ then gives $Y'\subseteq G$.
The same argument starts a supported component from any local ordinary
supported prime of sufficiently large quotient dimension.
These are the chart comparisons called
\srcid{ChartOfFlat} and \srcid{GlobOfCharts} in the source.

\subsubsection{An explicit connected family in characteristic $\ell$}
We prove the interpolation used for the characteristic-$\ell$ core.
Let $k$ be a field of odd characteristic, $\Gamma$ a profinite group,
$\delta:\Gamma\to\{\pm1\}$ a continuous surjection, and
$c\in\Gamma$ an involution with $\delta(c)=-1$.
Put $\Gamma_0=\ker\delta$.
Let $\phi:\Gamma\to k^\times$ be a locally constant character with
$\phi(c)=1$. On $\Gamma_0$ put $g^*=cg^{-1}c^{-1}$ and
$P(g)=\phi(g)$, so $P(g^*)=P(g)^{-1}$.
Fix $g_0\in\Gamma_0$ with $P(g_0)^2\ne1$ and local homomorphisms
$\Gamma_v\to\Gamma_0$, on each of which $P^2$ is nontrivial.

For a $k$-algebra $Q$, call a locally constant function
$b:\Gamma_0\to Q$ a strict linear datum if
\[
 \begin{split}
 b(gh)&=P(h)b(g)+P(g)^{-1}b(h),\\
 b(g^*)&=b(g),\\
 b(s)&=(P(s)-P(s)^{-1})\gamma_v
                \quad(s\in\Gamma_v)
 \end{split}
\]
for some $\gamma_v\in Q$.
Call a pair of locally constant functions $(y,a)$ a lower datum if
\begin{align*}
 y(gh)&=P(h)y(g)+y(h),\\
 a(gh)&=P(h)a(g)+y(g^*)y(h)+P(g)^{-1}a(h),\\
 a(g^*)&=a(g),\qquad y(g_0)=y(g_0^*),
\end{align*}
and, on each local group, there are $\alpha_v,\beta_v,\gamma_v\in Q$
with
\[
 \begin{split}
 y(s)&=(P(s)-1)\alpha_v,\\
 y(s^*)&=(1-P(s)^{-1})\beta_v,\\
 a(s)&=(P(s)-P(s)^{-1})\gamma_v
                  -(1-P(s)^{-1})\alpha_v\beta_v.
 \end{split}
\]
These formulas give the normalized lower-triangular representation
\[
 R_{y,a}(g)=
 \begin{pmatrix}
 P(g)&0&0\\
 y(g)&1&0\\
 a(g)&y(g^*)&P(g)^{-1}
 \end{pmatrix}.
\]
The two multiplication identities are precisely the matrix
multiplication identity. The involution gives
$R_{y,a}(g^*)=J\,{}^tR_{y,a}(g)J$ for the reversal matrix $J$,
and the equality at $g_0$ is the marked off-diagonal normalization.
The displayed local formulas imply the ordered three-character
ordinary determinant identity by multiplying the three triangular
factors. This checks the actual normalized representation conditions.

\begin{srlemma}\label{lem:sr-low-interpolation}
The lower data are stable under scalar extension and under
\[
             (y,a)\longmapsto(ty,t^2a+b),
\]
where $t\in Q$ and $b$ is a strict linear datum.
Suppose there exists a nonzero strict linear datum over $k$ and a lower
datum $(y_*,a_*)$ over $k$ with $y_*\ne0$.
Then the nonzero lower-data core is connected, and its subset
$y\ne0$ is dense, in every representing characteristic-$\ell$ core
having the field-point and polynomial-point universal properties
for these formulas.
\end{srlemma}
\begin{proof}
The operation preserves the multiplication and involution identities
by direct substitution. In the local formulas replace
\[
  (\alpha_v,\beta_v,\gamma_v)
       \quad\hbox{by}\quad
  (t\alpha_v,t\beta_v,t^2\gamma_v+\eta_v),
\]
where $b(s)=(P(s)-P(s)^{-1})\eta_v$.

Fix a nonzero strict linear datum $b_0$.
The pair $(0,b_0)$ is a lower datum.
For any field-valued lower datum $(y,a)$, the polynomial-valued pair
\[
            (Ty,\ T^2a+(1-T)b_0)
\]
is a lower datum over $Q[T]$. Its evaluations at one and zero are
$(y,a)$ and $(0,b_0)$.
The generic prime of this polynomial family specializes to both
endpoint primes. It stays in the nonzero core: a zero polynomial
datum would have zero evaluation at zero, contradicting $b_0\ne0$.
Thus every point is joined to the fixed point $(0,b_0)$ by a connected
three-point subset consisting of the two endpoint primes and the
generic prime of the family. Their union is connected.

Now take an exceptional point $(0,a)$.
The defining identities show that $a$ is a strict linear datum.
The family
\[
                    (Ty_*,\ T^2a_*+a)
\]
specializes to $(0,a)$ at zero, and its generic point has nonzero
first coordinate. Hence every exceptional point is in the closure
of the nonexceptional locus. Points already outside the exceptional
locus require no argument.
\end{proof}

The two nonzero data have a direct cohomological origin.
For $m=1,2$, let $\xi_m=\phi^{-m}\delta$, and suppose that the strict
degree-one cohomology for $\xi_m$, with local coboundaries at the
chosen groups, is nonzero.
These are precisely the two nontriviality fields in the source frame.
The local nontriviality of $P^2$ ensures that the corresponding
mapping-cone cohomology injects into $H^1(\Gamma,k(\xi_m))$:
if a global cocycle is a coboundary, the difference of every local
primitive and the global primitive is locally invariant, hence zero.
Restriction to $\Gamma_0$ is injective by averaging over the quotient
of order two.

For $m=2$, choose a nonzero class and a cocycle $x$ with $x(c)=0$,
which is possible because $\xi_2(c)=-1$.
Then
\[
                         b_0(g)=P(g)x(g)
\]
is a nonzero strict linear datum.
The cocycle equation proves its multiplication formula.
Normalization at $c$ gives
$x(g^*)=P(g)^2x(g)$, proving its involution formula.
A local coboundary $x(s)=(P(s)^{-2}-1)h_v$ gives the local formula
with $\gamma_v=-h_v$.

For $m=1$, the same normalization gives a nonzero locally constant
function $u(g)=P(g)x(g)$ with
\[
 u(gh)=P(h)u(g)+u(h),\qquad
 u(g^*)=P(g)^{-1}u(g),\qquad
 u(s)=(P(s)-1)\alpha_v.
\]
Put $p_0=P(g_0)$ and choose
\[
                 \lambda=-\frac{u(g_0)}{p_0+1}.
\]
The denominator is nonzero because $p_0^2\ne1$.
Define
\[
 \begin{split}
 y_*(g)&=u(g)+\lambda(P(g)-1),\\
 a_*(g)&=\lambda^2
   -\frac{\lambda^2}{2}(P(g)+P(g)^{-1})
   -\lambda(u(g)+P(g)^{-1}u(g))
   +\frac{P(g)^{-1}u(g)^2}{2}.
 \end{split}
\]
Substitution of the three identities for $u$ proves the lower-data
identities. At $g_0$ the difference
$y_*(g_0)-y_*(g_0^*)$ is
\[
       \frac{p_0-1}{p_0}
                  \bigl(u(g_0)+\lambda(p_0+1)\bigr)=0.
\]
For the local formulas use
\[
 \alpha_v'=\lambda+\alpha_v,\quad
 \beta_v'=\alpha_v-\lambda,\quad
 \gamma_v'=\frac{\alpha_v^2}{2}
                         -\lambda\alpha_v-\frac{\lambda^2}{2}.
\]
Finally $y_*=0$ would give $u=-\lambda(P-1)$.
Comparison with the involution formula at $g_0$ forces
$2\lambda(p_0-1)=0$, hence $\lambda=0$ and $u=0$, a contradiction.
This proves the required nonzero lower datum.

The source formulas and their universal-property application are in
\href{https://github.com/ChristianKudeba/alpoge-rh-verification/blob/alpoge-source-7cbf7a0c/src/Development/AutomorphicArithmetic/AutomorphicVarieties/FrobeniusSmoothVectors/StrSMul.lean}
{\srcid{StrSMul}},
\href{https://github.com/ChristianKudeba/alpoge-rh-verification/blob/alpoge-source-7cbf7a0c/src/Development/AutomorphicArithmetic/AutomorphicVarieties/FrobeniusSmoothVectors/ExistsNonzeroLowCocycle.lean}
{\srcid{ExistsNonzeroLowCocycle}}, and
\href{https://github.com/ChristianKudeba/alpoge-rh-verification/blob/alpoge-source-7cbf7a0c/src/Development/AutomorphicArithmetic/AutomorphicVarieties/FrobeniusSmoothVectors/ExistsNonzeroLowCocycle.lean}
{\srcid{preconnected_rank_three}}.
These arguments produce connectedness and density once the source
field-point and polynomial-point identifications have been established;
those identifications must not be inferred merely from the existence
of a formal local universal ring.

\subsubsection{The remaining specialized interface}
For the source's rank-three universal framed space $S$, the generic
and exceptional strata are
\[
 T=\{P:\mathfrak t_+\not\subseteq P,\
          \mathfrak t_-\not\subseteq P,\
          \mathfrak e\not\subseteq P\},\qquad
 E=\{P:\mathfrak t_+\not\subseteq P,\
          \mathfrak t_-\not\subseteq P,\
          \mathfrak e\subseteq P\},
\]
where $\mathfrak t_\pm$ are generated by the two triangular directions,
and $\mathfrak e$ by entries of odd weight for weights $(1,0,-1)$.
The supported locus is the zero locus of the Hecke ideal pulled back
to the framed entry algebra.
A selected prime has quotient dimension one, contains $\ell$, lies
off both triangular loci, and has pairwise separated local characters.

The local propagation threshold is $3d+1$ in rank three and $2d+1$
in rank two. The preceding lemmas prove the abstract propagation
once the following specialized facts are supplied:
\begin{enumerate}
\item The field-point and polynomial-point descriptions of the
characteristic-$\ell$ core, the flat local charts, and the dimension
bounds on the triangular and coincident-character loci.
\item The sharp local ideal: outside $E$ it is zero. At a point of $E$
the source constructs an even/odd presentation
\[
 R=\mathcal O'[[x_1,\ldots,x_a,y_1,\ldots,y_g]]
                         /(F_1,\ldots,F_b,G_1,\ldots,G_g),
 \qquad a=b+3d+1,\qquad G=M y .
\]
The even equations are $F_i$, the odd equations are $G_i$, and
$\mathfrak e$ is generated by the images of the $y_j$.
The sharp ideal is generated by the image of $\det M$.
The adjugate identity $(\operatorname{adj}M)G=(\det M)y$ shows directly
that every prime outside $E$ contains this ideal.
The required connectedness of its minimal-prime graph and the
ordinary/thick dichotomy still require the determinant-section and
deformation-obstruction arguments of
\srcid{Detsec}, \srcid{SumLe}, \srcid{ChainedOfComap}, and
\srcid{ThickOfDer}.
\item The Hecke-support implications P1 and P23, corresponding exactly
to hypotheses 6 and 7 of Lemma~\ref{lem:sr-local-propagation}, and a
seed of full weight-space dimension. The two-dimensional propagation
must also supply the finite exceptional support ideal used in rank
three. These are the inputs assembled in \srcid{PrLinkedEnd} and
\srcid{AbsMain}.
\end{enumerate}
Thus the prime avoidance, interpolation and propagation arguments
above are proved statements. They do not assert the full arithmetic
endpoint before these specialized facts have been proved.
In particular a theorem assuming that every residual constituent
is individually polarized cannot supply this interface merely from
the paired outer constituents $\bar\Psi$ and $\bar\Psi^{-1}$.

\subsection{One patched module and its completed supports}
\label{sec:sr-common-patched-module}
This section proves the common-module and completed-support argument
used in the source's \srcid{Raw} construction. In particular, the
modules at different points of a prime chain arise from one module.
Choosing a new module independently at each point would not prove
the support-transport assertion.

\subsubsection{The exact finite-level hypotheses}
Let $\Lambda$ be a complete Noetherian local ring with finite residue
field, and let $(R,\mathfrak m_R)$ be a Noetherian local ring.
Fix a homomorphism
\[
       s:\Lambda[[Y_1,\ldots,Y_r]]\longrightarrow R.
\]
Write $y_i=s(Y_i)$, $Y=(y_1,\ldots,y_r)$ and
$\mathfrak m=\mathfrak m_\Lambda R$.
Let $f\in R$, let $I\subseteq R$ be an ideal, and let $X$ be an
$R$-module finite over $\Lambda$, with compatible scalar actions and
$YX=0$. Fix a nonprincipal ultrafilter $\mathcal U$ on $\mathbb N$.
Let $M_N$ be $R$-modules satisfying:
\begin{enumerate}
\item For every $k,j\ge1$, for $\mathcal U$-almost all $N$ the module
\[
                  M_N/(Y^k+\mathfrak m_\Lambda^j)M_N
\]
is finite free over
$\Lambda[[Y_1,\ldots,Y_r]]/((Y_1,\ldots,Y_r)^k,\mathfrak m_\Lambda^j)$.
No common bound on the ranks is assumed.
\item For every $j\ge1$, for $\mathcal U$-almost all $N$ there are
$R$-linear maps
\[
 X/\mathfrak m_\Lambda^jX
       \ \mathop{\rightleftarrows}^{\alpha_{N,j}}_{\beta_{N,j}}\
             M_N/(Y+\mathfrak m_\Lambda^j)M_N
\]
such that both composites are multiplication by $f$.
\item There are an integer $e\ge1$ and a set of generators
$\mathscr R$ of $I$ such that, for every $a\in\mathscr R$ and $c\ge1$,
for $\mathcal U$-almost all $N$ there is $b\in\operatorname{Ann}_R M_N$
with
\[
                            (fa)^e-b\in\mathfrak m_R^c.
\]
The exponent $e$ precedes the approximation depth, the level, and
the relation generator.
\end{enumerate}
These hypotheses separate finite-level automorphic control from
the commutative algebra proved below. In the application,
$\Lambda$ is the fixed ordinary weight algebra with its extra scalar
variable, the $Y_i$ are the Taylor--Wiles variables, $f$ is the
auxiliary eigenvalue discriminant, and $I$ contains the kernel,
polarization, ordinary, inertia and two auxiliary relations.

\begin{srproposition}\label{prop:sr-one-patched-module}
Put $S=R[1/f]$ and $C=\widehat S^{\,Y}$, the $Y$-adic completion.
Under the three finite-level hypotheses there is one finite
$C$-module $\mathcal M$ with
\begin{equation}\label{eq:sr-patched-fiber}
 \mathcal M/Y\mathcal M\simeq C\otimes_R X,\qquad
 I C\subseteq\sqrt{\operatorname{Ann}_C\mathcal M}.
\end{equation}
Moreover
\begin{equation}\label{eq:sr-patched-gr}
 \operatorname{gr}_Y\mathcal M
           \simeq(\mathcal M/Y\mathcal M)[T_1,\ldots,T_r].
\end{equation}
Thus the variables are a quasi-regular sequence on $\mathcal M$.
\end{srproposition}
\begin{proof}
For $k\ge0$, define the $R$-module
\[
 U_k=\prod_N M_N\big/
 \left\{(z_N):
    \text{for every }j,\ 
     z_N\in(Y^k+\mathfrak m_\Lambda^j)M_N
                  \text{ for }\mathcal U\text{-almost all }N\right\}.
\]
In particular $U_0=0$.
The identity on representatives gives surjective maps
$U_{k+1}\to U_k$. Their kernels are $Y^kU_{k+1}$.

Here is the needed diagonal argument for the last assertion.
Choose generators $a_1,\ldots,a_t$ of $Y^k$.
If a representative $(z_N)$ maps to zero in $U_k$, then for each
$j$ it admits, for almost all $N$, a decomposition
\[
       z_N=\sum_{\nu=1}^t a_\nu w_{\nu,N}+v_N,
                          \qquad v_N\in\mathfrak m_\Lambda^jM_N.
\]
Choose a function $j(N)\to\infty$ along $\mathcal U$ on which these
decompositions hold. Such a choice is obtained by imposing the first
$j$ conditions and $N\ge j$, and taking the largest admissible
$j\le N$. Use a decomposition at $j(N)$ for each $N$.
The error sequence represents zero in $U_{k+1}$, which proves that
the kernel is generated by the $a_\nu$.
The reverse inclusion is immediate.
The same argument shows
\[
 U_k/Y^aU_k\simeq U_a\quad(1\le a\le k),\qquad Y^kU_k=0.
\]

We next identify $U_1$ after inverting $f$.
Choose a function $j(N)\to\infty$ along $\mathcal U$ and control maps
from hypothesis 2 at precision $j(N)$.
Lifting the images of elements of $X$ gives an $R$-linear map
$\alpha:X\to U_1$; all failures of linearity of arbitrary lifts vanish
in the defining quotient.
For a representative $(z_N)$ and fixed $j$, reduce
$\beta_{N,j(N)}(z_N)$ modulo $\mathfrak m_\Lambda^jX$ when
$j(N)\ge j$. The finite set $X/\mathfrak m_\Lambda^jX$ has a unique
$\mathcal U$-limit. These limits are compatible in $j$.
Since the finite $\Lambda$-module $X$ is complete, they define an
$R$-linear map $\beta:U_1\to X$.
Changing a representative by a null sequence does not change the limit:
$\beta_{N,j(N)}$ kills the $Y$-part and preserves the
$\mathfrak m_\Lambda^j$-part.
At each fixed precision the control identities give
\[
                      \alpha\beta=f,\qquad\beta\alpha=f.
\]
Consequently
\[
                         S\otimes_R U_1\simeq S\otimes_R X.
\]

Put $P_k=S\otimes_R U_k$.
Localization is exact, so $P_{k+1}\to P_k$ is surjective with kernel
$Y^kP_{k+1}$, and $Y^kP_k=0$.
Also $P_1$ is finite over $S$.
The inverse limit
$\mathcal M=\varprojlim_kP_k$ is therefore finite over $C$,
and $\mathcal M/Y^k\mathcal M\simeq P_k$ for every $k$.
We give the finite-generation detail to make this use of completion
explicit. Lift a finite generating set of $P_1$ compatibly through the
surjective tower. If $K_k$ is the kernel of the resulting surjection
$S^t\to P_k$, then
\[
             Y^kS^t\subseteq K_k,\qquad
             K_k=K_{k+1}+Y^kS^t.
\]
Successively correcting coefficient vectors makes their differences
belong to $Y^kS^t$; they converge in $C^t$.
This proves $C^t\to\varprojlim P_k$ is onto.
The same correction applied to relation vectors identifies the
kernel of the $k$-th projection with $Y^k\mathcal M$.
It also proves separatedness.

For completeness, hypothesis 1 gives the associated graded
calculation. A degree-$a$ relation
\[
               \sum_{|\nu|=a}y^\nu c_\nu
                           \in(Y^{a+1}+\mathfrak m_\Lambda^j)M_N
\]
in a module free over the truncated power-series ring forces
$c_\nu\in(Y+\mathfrak m_\Lambda^j)M_N$ for every $\nu$:
in each free coordinate the degree-$a$ monomials are independent over
$\Lambda/\mathfrak m_\Lambda^j$.
Only finitely many monomials occur.
Pass to $\mathcal U$-classes and use the diagonal kernel argument.
For $k>a$ this proves the same coefficient implication in $U_k$.
It survives localization and passage to the separated inverse limit.
The usual map from the polynomial module on $\mathcal M/Y\mathcal M$
to the associated graded is always surjective, and this coefficient
implication proves injectivity. This gives
\eqref{eq:sr-patched-gr}.
The identification for $k=1$ gives the first assertion of
\eqref{eq:sr-patched-fiber}.

Finally fix $a\in\mathscr R$ and $k,j\ge1$.
The finite $R$-module $X/\mathfrak m_\Lambda^jX$ has finite length,
so some power $\mathfrak m_R^c$ kills it.
The control maps at precision $j$ therefore give, for almost all $N$,
\[
                   f\mathfrak m_R^cM_N
                            \subseteq(Y+\mathfrak m_\Lambda^j)M_N.
\]
Taking the $k$-fold product of the acting ideals gives
\[
              f^k\mathfrak m_R^{ck}M_N
                            \subseteq(Y^k+\mathfrak m_\Lambda^j)M_N.
\]
Use hypothesis 3 at precision $ck$ and subtract an actual annihilator.
It follows that
\[
              f^{k+e}a^eM_N
                            \subseteq(Y^k+\mathfrak m_\Lambda^j)M_N
\]
for almost all $N$. The power $f^{k+e}$ is independent of $j$.
Hence $f^{k+e}a^e$ kills $U_k$.
After inverting $f$, $a^e$ kills $P_k$ for every $k$.
Separatedness proves $a^e\mathcal M=0$.
The generators of $I$ therefore map to the radical of the annihilator,
which completes the proof.
\end{proof}

This reconstructs the common module produced in
\srcid{AnnihilatorBaseChange.PmodEx} from the actual finite-level
freeness, paired control maps, and uniform nilpotence hypotheses.
It does not posit its existence separately at each specialization.

\subsubsection{Dimension at a completed point}
\begin{srlemma}\label{lem:sr-patched-support-dimension}
Let $L$ be a complete regular local ring of dimension $d$, let $B$ be
a complete Noetherian local $L$-algebra with local structure map,
and let $M$ be a finite nonzero $B$-module.
Suppose $y_1,\ldots,y_r\in\mathfrak m_B$ satisfy
\[
 \operatorname{gr}_{(y)}M
                  \simeq(M/(y)M)[T_1,\ldots,T_r],
\]
and $M/(y)M$ is finite free over $L$.
Then every irreducible component of $\operatorname{Supp}_BM$ has
dimension exactly $d+r$.
\end{srlemma}
\begin{proof}
Evaluation gives a local map $A=L[[Y_1,\ldots,Y_r]]\to B$ with
$Y_i\mapsto y_i$.
Lift an $L$-basis of $M/(y)M$ to $M$.
Successive approximation in powers of $(y)$ gives a surjection
$A^t\to M$: at each stage choose the next homogeneous coefficient,
and use completeness of the finite $B$-module to take the limit.
The associated graded hypothesis proves injectivity:
a nonzero relation has a least total degree in the $Y_i$, whose
coefficient vector would give a nonzero relation in the corresponding
free homogeneous piece.
Thus $M$ is finite free over $A$.

The faithful action of $B/\operatorname{Ann}_BM$ embeds that ring in
$\operatorname{End}_A(M)$, a finite free $A$-module.
Consequently $B/\operatorname{Ann}_BM$ is a finite torsion-free
$A$-algebra.
Every minimal prime contracts to zero in the domain $A$, since a
nonzero element of $A$ acts injectively on this torsion-free algebra
and thus lies in no minimal associated prime.
For each such prime $P$, the domain
$(B/\operatorname{Ann}_BM)/P$ is finite integral over $A$.
Its dimension is $\dim A=d+r$, proving the claim.
\end{proof}
This is the support-dimension content of
\srcid{AnnihilatorBaseChange.PmDimQuot}
and \srcid{AlgebraTools.CaCmPatch}.

In the application, the regular local ring $L$ is the completion of
the ordinary weight algebra at the contracted point.
The relevant special fiber is finite free over $L$.
Here is why finite freeness is preserved in that step.
The action of $R$ on the finite free $\Lambda$-module $X$ factors
through the finite $\Lambda$-algebra
$R/\operatorname{Ann}_RX\subseteq\operatorname{End}_\Lambda X$.
After completing $\Lambda$ at the contracted prime, this finite algebra
is a product of its complete local factors.
Completion at the selected $R$-prime picks one factor.
The corresponding piece of $X\otimes_\Lambda L$ is an idempotent
direct summand of a finite free $L$-module, and is therefore finite
free because $L$ is local. Zero summands have empty support.

\subsubsection{Fiber support and transport between completions}
Let $q_i$ be any finite family of primes of $R$ satisfying
\[
                        I+Y\subseteq q_i,\qquad f\notin q_i.
\]
For such a prime, $q_iC$ is prime and contracts to $q_i$.
Indeed $C/YC\simeq R[1/f]/Y$, so primes containing $Y$ are determined
by the ordinary localization correspondence.
Completing $C_{q_iC}$ gives the same ring as completing $R_{q_i}$:
their quotients modulo every power of the selected maximal ideal
agree, because $Y\subseteq q_i$ and $f$ is a unit there.
Write this common completion as $B_i^0$, and put
\[
 \mathcal M_i=B_i^0\otimes_C\mathcal M,\quad
 B_i=B_i^0/IB_i^0,\quad
 S_i=(\operatorname{Ann}_{B_i^0}\mathcal M_i+IB_i^0)/IB_i^0.
\]
The radical assertion in \eqref{eq:sr-patched-fiber} ensures that
the minimal support components and their dimensions are unchanged
by this quotient.

\begin{srlemma}\label{lem:sr-completed-support}
The ideals $S_i$ satisfy:
\begin{enumerate}
\item If the completed special fiber is finite free over a regular
local weight ring of dimension $d$, then
$\dim B_i/P=d+r$ for every $P\in\operatorname{Min}(S_i)$.
\item For every prime $Q\subseteq B_i$ containing $Y$,
\[
         S_i\subseteq Q
          \quad\Longleftrightarrow\quad
          \operatorname{Ann}_RX\subseteq Q\cap R.
\]
\item Suppose $Z\subseteq Q$ are primes of $B_i$, $S_i\subseteq Z$,
$Y\subseteq Q$, and $Q\cap R\subseteq q_j$.
There is a prime $Z'\subseteq B_j$ with
\[
                       S_j\subseteq Z',
                  \qquad Z'\cap(R/I)=Z\cap(R/I).
\]
\end{enumerate}
\end{srlemma}
\begin{proof}
Flat base change to a completion preserves the associated graded
description and the special fiber. The first assertion therefore
follows from Lemma~\ref{lem:sr-patched-support-dimension}.

For the second, localize at the inverse image of $Q$ in $B_i^0$.
Since the $y_i$ are in that maximal ideal and $\mathcal M_i$ is finite,
Nakayama gives
\[
 (\mathcal M_i)_Q\ne0
        \quad\Longleftrightarrow\quad
      (\mathcal M_i/Y\mathcal M_i)_Q\ne0.
\]
Use \eqref{eq:sr-patched-fiber}.
For any ring map $A\to B$ and finite $A$-module $N$,
\[
     \operatorname{Supp}_B(B\otimes_AN)
              =(\operatorname{Spec}B\to\operatorname{Spec}A)^{-1}
                                      \operatorname{Supp}_AN.
\]
To verify this identity, tensor at a prime with its residue field.
The resulting tensor product is the extension of
$N\otimes_A\kappa(p)$ through the field inclusion
$\kappa(p)\hookrightarrow\kappa(Q)$, so it is nonzero exactly when
the former vector space is nonzero.
For finite modules, support is the zero locus of the annihilator.
This proves assertion 2.

For assertion 3, first work before quotienting by $I$.
Let $p=Z\cap C$.
The prime $p$ lies in $\operatorname{Supp}_C\mathcal M$, by the
same support identity.
Since $Q$ contains $Y$, its contraction to $C$ is determined by
its contraction to $R$ modulo $Y$; hence $Q\cap R\subseteq q_j$
implies $Q\cap C\subseteq q_jC$.
Thus $p\subseteq q_jC$.
The completion of $C_{q_jC}$ is faithfully flat, so there is a prime
$Z'$ above $p$. The support identity puts it in the support of the
base change of the same module $\mathcal M$.
Contraction to $R$ is unchanged.
The ideal $I$ lies in every support prime by the radical annihilation
assertion, so the primes descend to the displayed quotients.
This proves assertion 3.
\end{proof}
The source separates these three facts as
\srcid{PmDimQuot}, \srcid{PmFib}, and \srcid{PmMove}.
The last proof explicitly uses the common module: equality of numerical
dimensions of independently chosen modules would not give transport
of a support prime.

\subsubsection{The completed-support chain}
Here is the final algebraic form of the source's \srcid{ChainAbs}.
Let $q_0,\ldots,q_h$ be selected points as above.
Assume that the completed support components have dimension at least
$d_0$. Let $\mathcal T$ be a property of primes of $R/I$ stable under
generalization, and let $\mathcal E_i$ be a property of primes of $B_i$.
Suppose the following local inclusion statement holds:
\[
 \left.
 \begin{array}{c}
 Z,Q\in\operatorname{Spec}B_i,\quad
 \mathcal T(Z\cap(R/I)),\quad\dim B_i/Z\ge d_0,\\
 \mathcal T(Q\cap(R/I))\ \text{or}\ \mathcal E_i(Q)
 \end{array}\right\}
                           \ \Longrightarrow\ Z\subseteq Q.       \tag{*}
\]
Suppose bridge primes $E_i$ of $R/I$ contain $Y$, lie below both
$q_i/I$ and $q_{i+1}/I$, and every prime of $B_i$ above $E_i$ has
property $\mathcal E_i$.
Suppose there is a prime $p_0\subseteq q_0/I$ containing $Y$, with
$\mathcal T(p_0)$ and
$\operatorname{Ann}_RX\subseteq p_0\cap R$.

\begin{srproposition}\label{prop:sr-completed-support-chain}
If $p\subseteq q_h/I$ contains $Y$ and satisfies either
$\mathcal T(p)$ or
$\mathcal E_h(Q)$ for every prime $Q$ of $B_h$ above $p$, then
\[
                            \operatorname{Ann}_RX\subseteq p\cap R.
\]
\end{srproposition}
\begin{proof}
Faithful flatness lifts $p_0$ to a prime of $B_0$.
The fiber-support equivalence puts it in $V(S_0)$.
Choose a minimal support prime $Z_0$ below it.
Then $\mathcal T(Z_0\cap(R/I))$ holds by generalization, and its
dimension is at least $d_0$.

Inductively suppose there is a supported prime $Z_i$ of $B_i$
whose contraction has property $\mathcal T$.
Replace it by a minimal support prime underneath.
Lift the bridge $E_i$ to a prime $Q_i$ of $B_i$.
The bridge hypothesis gives $\mathcal E_i(Q_i)$, so (*) gives
$Z_i\subseteq Q_i$.
The support-transport lemma then moves $Z_i$ to a supported prime
of $B_{i+1}$ with the same contraction to $R/I$.
This preserves $\mathcal T$ and proves the induction.

At the last point, choose a minimal support prime beneath the
supported prime obtained by induction, and lift $p$ to $Q\in B_h$.
The final hypothesis and (*) put the minimal support prime below $Q$.
Hence $S_h\subseteq Q$.
The fiber-support equivalence now gives
$\operatorname{Ann}_RX\subseteq p\cap R$, as required.
\end{proof}

\paragraph{What this closes in the actual source route.}
The common-module construction, its uniform nilpotence passage,
its completed support dimensions, the fiber identity, support
transport, and the chain argument have now been proved from the
displayed finite-level hypotheses and (*).
The source modules are
\href{https://github.com/ChristianKudeba/alpoge-rh-verification/blob/alpoge-source-7cbf7a0c/src/Development/AutomorphicArithmetic/AutomorphicVarieties/PatchingModules/AnnihilatorBaseChange/PmodEx.lean}
{\srcid{PmodEx}},
\href{https://github.com/ChristianKudeba/alpoge-rh-verification/blob/alpoge-source-7cbf7a0c/src/Development/AutomorphicArithmetic/AutomorphicVarieties/PatchingModules/AutomorphicDetection/AnnBc.lean}
{\srcid{AnnBc}}, and
\href{https://github.com/ChristianKudeba/alpoge-rh-verification/blob/alpoge-source-7cbf7a0c/src/Development/AutomorphicArithmetic/AutomorphicVarieties/PatchingModules/AnnihilatorBaseChange/ChainAbs.lean}
{\srcid{ChainAbs}}.
The single universal ring and relation ideal are defined in
\srcid{SpecializationMapCompUniversalSection}; \srcid{LiftPt.core}
applies the preceding facts to them.
The arithmetic local inclusion (*) is the separate producer
\srcid{AutomorphicDetection.GalOut}. It uses \srcid{S6} to realize the
completed determinant over a finite quadratic coefficient algebra,
then the first- and second-Selmer bounds and the Laurent presentation
\srcid{SelmerScaling.PresOut0}.
The finite-level freeness, paired control maps, uniform relation
nilpotence, and that Galois inclusion have to be verified in this
same universal ring before the full \srcid{Raw} endpoint is asserted.

\subsection{The determinant algebra at paired components}
\label{sec:sr-paired-determinant-algebra}
We prove the last commutative-algebra step of the local inclusion
used in the completed-support chain. It treats the two paired blocks
through a finite coefficient algebra with an odd summand.

\subsubsection{The generic incidence ideal}
Let $g\ge1$, let $R$ be a Noetherian domain, and let
\[
 P_0=R[Z_{ij},T_j:1\le i,j\le g],\qquad
 J=((ZT)_1,\ldots,(ZT)_g,\det Z).
\]
\begin{srlemma}\label{lem:sr-generic-incidence}
The radical of $J$ is prime and contains no nonzero element of $R$.
It does not contain $T_1$. The same assertions hold after extending
the homogeneous ideal to $R[[Z_{ij},T_j]]$.
If the latter power-series ring is Noetherian, the height of this
prime is at most $g$.
\end{srlemma}
\begin{proof}
Put $E=\operatorname{diag}(0,1,\ldots,1)$.
Use two generic $g$-by-$g$ matrices $N,U$ and a scalar $t$ and define
a polynomial map on coordinate rings by
\[
                 Z\longmapsto NE\operatorname{adj}U,\qquad
                 T\longmapsto tUe_1.
\]
The image satisfies $ZT=0$ and $\det Z=0$.
Its coordinate ring is a subring of the polynomial domain
$R[N_{ij},U_{ij},t]$, so the kernel $K$ is prime.
The images of a nonzero constant and of $T_1=tU_{11}$ are nonzero.

Every field-valued pair $(M,v)$ with $Mv=0$ and $\det M=0$
factors through this parametrization.
If $v\ne0$, take the first column of an invertible $U$ to be $v$
and take $t=1$.
If $v=0$, take its first column to be any nonzero vector in $\ker M$
and take $t=0$.
In either case $MU$ has first column zero.
Thus
\[
         N=(\det U)^{-1}MU
       \quad\Longrightarrow\quad
         NE\operatorname{adj}U=M,\qquad tUe_1=v.
\]
Apply this factorization to the fraction field of every prime quotient
of $P_0$ containing $J$.
Every element of $K$ belongs to all such primes, so $K=\sqrt J$.

The ideal $J$ is homogeneous for total degree, and therefore so is $K$.
For a homogeneous finitely generated ideal $K\subseteq R[\mathbf Z]$,
a formal series belongs to $KR[[\mathbf Z]]$ exactly when each of its
homogeneous polynomial pieces belongs to $K$.
Indeed choose finite homogeneous generators and sum homogeneous
coefficient decompositions degree by degree.
If two series are outside this extended ideal, choose the least
degrees $a,b$ of homogeneous pieces outside $K$.
The degree-$a+b$ piece of their product is, modulo $K$, the product
of those two pieces, and is nonzero because $K$ is prime.
Thus the extended ideal is prime.
Its contraction is $K$, and it is the radical of the extension of $J$.

Write $\widehat K$ for that prime.
Choose a prime $Q$ minimal over the $g$ row relations $ZT$
inside $\widehat K$.
Since $T_1\notin\widehat K$, also $T_1\notin Q$.
The adjugate identity gives $(\det Z)T_1\in Q$, hence $\det Z\in Q$.
Therefore $\widehat K=Q$.
Krull's height theorem gives $\operatorname{ht}\widehat K\le g$.
\end{proof}
This is the explicit parametrization and homogeneous-completion
argument of \srcid{CommutativeAlgebraHeight.GDetrad},
\srcid{GHomprime} and \srcid{GDetps}.

Now let $\mathcal O$ be a complete DVR and work in
\[
 P=\mathcal O[[Z_{ij},T_j,W_a]].
\]
The $g$ row relations $ZT$ form a regular sequence.
One direct check gives diagonal $Z_{ii}$ weight one,
off-diagonal $Z_{ij}$ weight two, and every $T_j$ weight one.
The initial forms of the rows are the monomials $Z_{ii}T_i$.
They involve disjoint sets of variables and form a regular sequence.
Taking the least weighted degree of a relation then proves regularity
of the original rows, one row at a time.
Since $P$ is regular local and catenary, Lemma~\ref{lem:sr-generic-incidence}
therefore gives
\begin{equation}\label{eq:sr-incidence-dimension}
               \dim P/\sqrt{(ZT,\det Z)}=\dim P-g.
\end{equation}

The same conclusion applies to elements $z_{ij},t_j$ of a power-series
ring over $\mathcal O$ whose classes are linearly independent modulo
$(\varpi,\mathfrak m_P^2)$.
Complete the independent classes to a relative coordinate basis.
The homomorphism from the power-series ring in those coordinates is
an isomorphism: successive approximation gives surjectivity, and its
invertible linear term gives an inverse by solving each homogeneous
degree. Constants in $\varpi\mathcal O$ cause no problem because the
map is continuous for the maximal-ideal topology.
Thus the displayed incidence ideal is carried to the generic one.

\subsubsection{Stability detected by odd products}
Let $(B,\mathfrak m_B)$ be local and let $A$ be a finite commutative
$B$-algebra with an injective structure map and a decomposition
\[
                          A=B\oplus N,\qquad NN\subseteq\mathfrak m_B.
\]
Define the ideal
\[
  I_{\rm odd}=(b\in B:\ uv=b\text{ in }A
                              \text{ for some }u,v\in N).
\]
Consider a representation whose $B$-linear span is exactly
\[
 \begin{pmatrix}
 M_a(B)&M_{a,n-a}(N)\\
 M_{n-a,a}(N)&M_{n-a}(B)
 \end{pmatrix},\qquad0<a<n .
\]
\begin{srlemma}\label{lem:sr-odd-product-detector}
At a prime $q$ of $B$, this representation has full matrix span over
an algebraic closure of $\kappa(q)$ if and only if
$I_{\rm odd}\not\subseteq q$.
\end{srlemma}
\begin{proof}
Pass to a prime of the integral algebra $A$ over $q$ and then to its
residue field.
If $I_{\rm odd}\subseteq q$, every $u\in N$ has square zero in that
field, hence has zero image. All off-block entries vanish, so the span
is block diagonal and is not the full matrix algebra.

Otherwise choose $u,v\in N$ with $uv=b\notin q$.
The image of $u$ in any residue field above $q$ is nonzero.
The given span contains every same-block matrix unit and every
off-block matrix unit multiplied by $u$.
After extending scalars, division by the nonzero image of $u$ gives
every matrix unit. The span is the full matrix algebra.
\end{proof}
For a single block ($a=n$), the given span is already the full matrix
algebra; in the source the odd module is then zero.
The trace-Gram determinant criterion gives the same statement without
choosing an algebraic closure: the trace pairing on the full matrix
algebra is nondegenerate, so a nonzero $n^2$-by-$n^2$ Gram minor is
equivalent to full span. This is exactly the determinant predicate
\srcid{DiscSt} and the proof in \srcid{ClosureConjugation.UrepStable}.

\subsubsection{The ordinary presentation}
\begin{srlemma}\label{lem:sr-ordinary-presentation-prime}
Let $\mathcal O\to B$ be a local homomorphism from a complete DVR to
a complete Noetherian local ring, inducing an isomorphism on residue
fields. Suppose $x_1,\ldots,x_D\in\mathfrak m_B$ satisfy
\[
            \mathfrak m_B\subseteq
                 (x_1,\ldots,x_D)+\varpi B+\mathfrak m_B^2.
\]
If $Z$ is prime and $\dim B/Z\ge D+1$, then $Z=0$.
Consequently $Z\subseteq Q$ for every prime $Q$ of $B$.
\end{srlemma}
\begin{proof}
The displayed cotangent generation and completeness give a surjection
$\pi:\mathcal O[[X_1,\ldots,X_D]]\twoheadrightarrow B$,
$X_i\mapsto x_i$, by successive approximation.
The source is a regular local domain of dimension $D+1$.
If its kernel after composition with $B\to B/Z$ contained a nonzero
element, the principal ideal theorem and the regular local dimension
formula would give $\dim B/Z\le D$.
Hence the composite has zero kernel.
Surjectivity now forces both $\ker\pi=0$ and $Z=0$.
\end{proof}
This proves the complete algebraic content of \srcid{PresOut0}.

\subsubsection{The paired presentation}
The following lemma is the last step in the exceptional local
inclusion. It states the exact presentation data, including the
independence that the Selmer calculation must supply.

\begin{srlemma}\label{lem:sr-paired-presentation-prime}
Let $\mathcal O\to B$ have the hypotheses of
Lemma~\ref{lem:sr-ordinary-presentation-prime}, and let
$A=B\oplus N$ be a finite commutative $B$-algebra as above.
Let $m_1,\ldots,m_g$ generate $N$ over $B$, and let
$x_1,\ldots,x_D\in\mathfrak m_B$.
Suppose the vectors
\[
                 \bar c_{ij}=(c_{ij,1},\ldots,c_{ij,D})\bmod\varpi
                    \quad(1\le i,j\le g)
\]
are linearly independent over the residue field, and put
\[
                  H=\varpi B+\mathfrak m_B^2+I_{\rm odd}.
\]
Assume
\[
 \mathfrak m_B\subseteq(x_1,\ldots,x_D)+H,\qquad
        \sum_j\left(\sum_a c_{ij,a}x_a\right)m_j\in HN
                                       \quad(1\le i\le g).
\]
If $Z,Q$ are primes of $B$ with
\[
 \dim B/Z\ge D+1,\qquad I_{\rm odd}\not\subseteq Z,\qquad
 I_{\rm odd}\not\subseteq Q\ \text{or}\ \operatorname{Ann}_BN\subseteq Q,
\]
then $Z\subseteq Q$.
\end{srlemma}
\begin{proof}
If $g=0$, then $N=0$ and $I_{\rm odd}=0$, contradicting the hypothesis
at $Z$. Thus $g\ge1$.
Since the $m_j$ generate $N$, choose $w_{ij}\in H$ such that
\[
                \sum_j\zeta_{ij}m_j=0,\qquad
          \zeta_{ij}=\sum_a c_{ij,a}x_a-w_{ij}.
\]
Multiplication by the adjugate proves
$(\det\zeta)m_j=0$ for every $j$, hence
\begin{equation}\label{eq:sr-det-ann-odd}
                             \det\zeta\in\operatorname{Ann}_BN.
\end{equation}
If $I_{\rm odd}\not\subseteq Q$, choose an odd product
$uv=b\notin Q$. Equation~\eqref{eq:sr-det-ann-odd} gives
$(\det\zeta)b=0$, using the injection $B\hookrightarrow A$.
Primality then gives $\det\zeta\in Q$.
If instead $\operatorname{Ann}_BN\subseteq Q$, the same conclusion
follows directly. Thus both alternatives at $Q$ give
$\det\zeta\in Q$.

We make the formal coordinates explicit.
The ring $A$ is complete local with maximal ideal
$\mathfrak m_BA+N$, since the image of $N$ modulo $\mathfrak m_BA$
has square zero. Every element of $I_{\rm odd}$ maps into
$\mathfrak m_A^2$.
The cotangent-generation assumption and the generators $m_j$
therefore give a surjection
\[
        \pi:P=\mathcal O[[X_1,\ldots,X_D,T_1,\ldots,T_g]]
                         \twoheadrightarrow A,
          \qquad X_a\mapsto x_a,\quad T_j\mapsto m_j .
\]
Lift each $w_{ij}$ to
$\widetilde w_{ij}\in(\varpi,\mathfrak m_P^2)$ and put
\[
                    z_{ij}=\sum_a c_{ij,a}X_a-\widetilde w_{ij}.
\]
The vectors $\bar c_{ij}$ are independent, so the $z_{ij}$ together
with the $T_j$ are independent in the relative cotangent space.
They extend to formal coordinates, as above.
The row relations $zT$ belong to $\ker\pi$.

Choose primes $Z',Q'$ of $A$ over $Z,Q$ by integrality.
At least one $m_j$ is outside $Z'$; otherwise all of $N$ lies in
$Z'$, so $I_{\rm odd}\subseteq Z'\cap B=Z$, a contradiction.
Put $P_Z=\pi^{-1}Z'$.
Then $zT\subseteq P_Z$ and at least one $T_j\notin P_Z$.
The adjugate identity gives $\det z\in P_Z$.
Consequently the prime
\[
                           K=\sqrt{(zT,\det z)}
\]
is contained in $P_Z$.
Lemma~\ref{lem:sr-generic-incidence} and
\eqref{eq:sr-incidence-dimension} give
\[
                        \dim P/K=\dim P-g=D+1.
\]
Since $A/Z'$ is finite integral over $B/Z$,
\[
                         \dim P/P_Z=\dim B/Z\ge D+1 .
\]
A strict inclusion $K\subsetneq P_Z$ in the complete regular local
ambient ring would lower the quotient dimension by at least one.
Hence $P_Z=K$.

Finally $P_Q=\pi^{-1}Q'$ contains $zT$ and $\det z$, because
$\det\zeta\in Q$. Thus
$P_Z=K\subseteq P_Q$, giving $Z'\subseteq Q'$.
Contracting to $B$ gives $Z\subseteq Q$.
\end{proof}
The proof supplies the power-series version needed for the source's
\srcid{CotangentIndependentOfResidues.presOut1}.
In particular it displays why the coefficient vectors must be
independent before choosing the formal coordinates.

\subsubsection{The actual local inclusion in the patched chain}
For the completed determinant algebra used in
Section~\ref{sec:sr-common-patched-module}, the source constructs
a finite commutative coefficient algebra $A=B\oplus N$ with
$NN\subseteq\mathfrak m_B$ and a matrix representation whose span is
the two-block algebra above.
The stable locus is detected by $I_{\rm odd}$ through
Lemma~\ref{lem:sr-odd-product-detector}.
The thick exceptional locus lies in
$V(\operatorname{Ann}_BN)$.

The first- and second-Selmer bounds and the spanning by the chosen
odd cocycles produce the presentation in
Lemma~\ref{lem:sr-paired-presentation-prime}:
\srcid{AlCSmul.pres} supplies the $x_a,m_j,c_{ij,a}$ with the
displayed generation, independence and relation properties.
In the single-block case $N=0$ it supplies
Lemma~\ref{lem:sr-ordinary-presentation-prime}.
Thus, once that finite coefficient algebra and those explicit Selmer
presentation data are constructed, the local implication (*) in
Proposition~\ref{prop:sr-completed-support-chain} is a proved
consequence of the two presentation lemmas.
The Galois source performs these two cases in
\srcid{AutomorphicDetection.GalOut.ctx_out}.

The algebra $A=B\oplus N$ is not an assumption that the two residual
constituents are separately polarized.
In the source it is constructed from the completed
Cayley--Hamilton algebra with its transpose involution, lifting its
orthogonal rank-one idempotents and identifying the two opposite
off-block corner modules with the odd summand $N$.
Its multiplication is
\[
       (b,u)(c,v)=(bc+\beta(u,v),\,bv+cu),
                    \qquad\beta:N\times N\longrightarrow B.
\]
The pairing is symmetric, and associativity of the corner products
gives $\beta(u,v)w=\beta(v,w)u$.
Residual block diagonality gives $\beta(N,N)\subseteq\mathfrak m_B$.
The construction, including determinant and involution compatibility,
is in \srcid{ClosureConjugation.OBLt},
\srcid{ComplOddPart}, \srcid{GrI} and \srcid{UrepOf},
assembled in \srcid{AutomorphicDetection.S6}.
Those constructions and the Selmer presentation are the remaining
arithmetic producers; the ordinary/paired prime-inclusion algebra
above has been proved independently.

\paragraph{Source locations.}
The generic incidence sources are
\href{https://github.com/ChristianKudeba/alpoge-rh-verification/blob/alpoge-source-7cbf7a0c/src/Development/AutomorphicArithmetic/AutomorphicVarieties/CommutativeAlgebraHeight/GDetrad.lean}
{\srcid{GDetrad}} and
\href{https://github.com/ChristianKudeba/alpoge-rh-verification/blob/alpoge-source-7cbf7a0c/src/Development/AutomorphicArithmetic/AutomorphicVarieties/CommutativeAlgebraHeight/GDetps.lean}
{\srcid{GDetps}}.
The completed local inclusion is
\href{https://github.com/ChristianKudeba/alpoge-rh-verification/blob/alpoge-source-7cbf7a0c/src/Development/AutomorphicArithmetic/AutomorphicVarieties/PatchingModules/AutomorphicDetection/GalOut.lean}
{\srcid{GalOut}}.
The exact paired presentation is
\href{https://github.com/ChristianKudeba/alpoge-rh-verification/blob/alpoge-source-7cbf7a0c/src/Development/AutomorphicArithmetic/AutomorphicVarieties/PatchingModules/SelmerScaling/CotangentIndependentOfResidues.lean}
{\srcid{CotangentIndependentOfResidues}},
and its Selmer producer is
\href{https://github.com/ChristianKudeba/alpoge-rh-verification/blob/alpoge-source-7cbf7a0c/src/Development/AutomorphicArithmetic/AutomorphicVarieties/PatchingModules/SelmerScaling/AlCSmul.lean}
{\srcid{AlCSmul}}.
